\section{Introduction}

Vol.~1 of this series established that time and space emerge from the causal log
$C_{ij}$ through self-referential closure dynamics~\cite{VED_vol1}.
Vol.~2 showed that the gauge group
$\mathrm{SU}(3)\times\mathrm{SU}(2)\times\mathrm{U}(1)$,
particle spectrum, and mass generation
can be interpreted as structural consequences of triangular closure~\cite{VED_vol2}.

The present volume addresses the remaining question:
how does gravitational dynamics arise within the same framework?

The starting point is the distance definition of Vol.~1:
\begin{equation}
  d_{ij} = -\log \tilde{S}_{ij},
  \qquad
  \tilde{S}_{ij} = \max_{\gamma: i\to j} \prod_{(ab)\in\gamma} S_{ab}.
\end{equation}
This definition is purely structural, determined by the most permeable path
in the causal log network.
It yields a consistent geometry, but does not produce gravity.
The reason is that it assumes structural homogeneity:
no mechanism biases path selection according to accumulated closure structure.

Gravity, within VED, appears precisely as such a bias.
The central claim of this volume is:
\begin{equation}
  \boxed{
    \text{Gravity is not a new force.
    It appears as a modification of the definition of distance.}
  }
\end{equation}

Incorporating closure density $L_{ab}$ into the path cost unifies geometry
and gravitational effect within a single selection principle.
From this modification, the Einstein equation emerges as the linear approximation
of a fixed-point structure, the Schwarzschild solution appears as an exterior limit,
the black hole interior is identified as a non-equilibrium closure phase,
and cosmic expansion arises as the global $\sigma > 0$ regime.

This interpretation is structural and qualitative in nature.
No quantitative predictions are derived, and no claims are made about
numerical correspondence with observed values.
The goal is to exhibit a coherent geometric reading of gravitational physics.

The paper is organized as follows.
Section~\ref{sec:gravity_distance} introduces the modified distance definition.
Section~\ref{sec:closure_density} develops the closure tensor and its relation
to the metric.
Section~\ref{sec:einstein} derives the Einstein equation as a fixed-point structure.
Section~\ref{sec:schwarzschild} examines the exterior solution.
Section~\ref{sec:black_hole} treats the black hole interior as a non-equilibrium phase.
Section~\ref{sec:cosmology} addresses cosmic expansion.
Section~\ref{sec:synthesis} synthesizes the three volumes.
Section~\ref{sec:discussion} discusses scope and limitations.
